% paper/sections/03b_levelset_mapping.tex
%
% §3 CLS 法の続き（03_levelset.tex からの分割）
% 内容：ψ-φ 関係・逆変換・曲率の数値計算と安定化
%
\subsection{\texorpdfstring{$\psi$}{ψ} と \texorpdfstring{$\phi$}{φ} の関係}
\label{sec:levelset_mapping}

\begin{tcolorbox}[defbox, title={$\psi$--$\phi$ 対応付けと解析的逆変換（ロジット関数）}]
Conservative Level Set 変数 $\psi$ と符号付き距離関数 $\phi$ は以下の関係にある：
\begin{equation}
  \psi(\bm{x},t) \equiv H_\varepsilon(\phi(\bm{x},t)), \qquad \bnabla\psi = \delta_\varepsilon(\phi)\,\bnabla\phi
\end{equation}
物性値補間 $\rho(\psi), \mu(\psi)$ には $\psi$ を直接用いる．
曲率 $\kappa$ は $\psi$ の微分値から直接計算できる
（定理~\ref{thm:curvature_invariance}，§~\ref{sec:curvature_direct_psi_cls}）．
Extension PDE（§~\ref{sec:extension_pde}）による界面越し場延長で $\phi$ が必要な場合は
逆変換（式~\eqref{eq:logit_inverse}）を用いる．

\medskip
\textbf{逆変換（$\psi \to \phi$）：}
CLS 法の移流計算は $\psi$ で行うが，幾何量計算には $\phi$ が必要である．
$H_\varepsilon(\phi) = 1/(1+e^{-\phi/\varepsilon})$ はシグモイド型関数であり，
その逆関数（ロジット関数）は\textbf{解析的に求まる}：
\begin{equation}
  \phi = H_\varepsilon^{-1}(\psi) = \varepsilon\ln\!\left(\frac{\psi}{1-\psi}\right)
  \label{eq:logit_inverse}
\end{equation}
$H_\varepsilon$ は単調増加関数（$H_\varepsilon' > 0$）であるから，解は一意に定まる．
唯一の実装上の注意は，$\psi \to 0$ または $\psi \to 1$ の\textbf{飽和領域}で
対数引数が $0$ または $\infty$ に近づき浮動小数点オーバーフローが生じる点である．
この場合は $|\phi| = \phi_{\max}$ に飽和させる（§~\ref{sec:newton_inversion} 参照）．
\end{tcolorbox}

\subsubsection{逆変換の実装：解析的逆関数と飽和処理}
\label{sec:newton_inversion}

式~\eqref{eq:logit_inverse} で示したように，逆変換は\textbf{ロジット関数}によって点ごとに $\Ord{1}$ で実行できる．
実装上の唯一の注意点は飽和領域における浮動小数点オーバーフローの回避である．

以下に実装の全手順を示す．

\begin{tcolorbox}[algbox, title={$\psi \to \phi$ 逆変換の実装手順（全格子点に対して点別に実行）}]
\begin{enumerate}
  \item \textbf{飽和判定：} $\psi_i < \delta_{\text{clamp}}$ または $\psi_i > 1 - \delta_{\text{clamp}}$（例：$\delta_{\text{clamp}} = 10^{-6}$）を確認する．
  \item \textbf{飽和領域：} 飽和している点では $\phi_i = \mp\phi_{\max}$ に固定する．
        推奨値は $\phi_{\max} = \varepsilon\ln\!\bigl((1-\delta_{\text{clamp}})/\delta_{\text{clamp}}\bigr) \approx 13.8\varepsilon$（$\delta_{\text{clamp}}=10^{-6}$ のとき）．
  \item \textbf{標準領域：} 飽和していない点では解析的逆変換
        \[
          \phi_i = \varepsilon\ln\!\left(\frac{\psi_i}{1-\psi_i}\right)
        \]
        を直接計算する．計算コストは対数1回のみ（点あたり $\Ord{1}$）．
\end{enumerate}

\textbf{備考：}
\begin{itemize}
  \item 飽和域を除き，ニュートン反復は\textbf{不要}である．解析的逆関数は残差ゼロを代数的に保証する．
  \item 飽和処理後，逆変換された $\phi$ のプロファイルが $|\bnabla\phi|\approx 1$ を満たしているかを
        確認するため，CLS 再初期化方程式（式~\eqref{eq:cls_reinit_iso}）を数回適用することが推奨される．
  \item 飽和領域の格子点は界面から遠く（$|\phi| \gg \varepsilon$），
        $|\bnabla\psi| \approx 0$ であるため曲率計算に寄与しない．
\end{itemize}
\end{tcolorbox}

飽和領域（$\psi_i < \delta_{\text{clamp}}$ または $\psi_i > 1-\delta_{\text{clamp}}$）で解析的逆変換が浮動小数点上限に近づく場合のフォールバックとして，以下の Newton 更新式を用いる（通常の格子点では不要；$\delta_\text{clamp}=10^{-6}$ 設定時，飽和域の対象格子点数は実質ゼロ）：
\begin{equation}
  \phi^{(k+1)} = \phi^{(k)} - \frac{H_\varepsilon(\phi^{(k)}) - \psi_i}{\dfrac{1}{\varepsilon}H_\varepsilon(\phi^{(k)})\bigl(1-H_\varepsilon(\phi^{(k)})\bigr)}
  \label{eq:newton_inversion}
\end{equation}
証明は飽和処理 $|\phi| \le \phi_{\max}$ によって反復が有界閉区間に閉じ込められる場合に成立する．通常の実装では飽和クランプがこれを自動的に保証する（付録~\ref{app:newton_quad_conv} 参照）．


\begin{tcolorbox}[mybox, title={ロジット逆変換 vs. Sussman 再初期化：$\psi \to \phi$ 変換の選択}]
本節で確立したロジット解析逆変換（式~\eqref{eq:logit_inverse}）と
標準 LS 法で使われる Sussman 再初期化を比較する．
\begin{center}
\renewcommand{\arraystretch}{1.3}
\begin{tabular}{lll}
\hline
 & \textbf{ロジット逆変換（本稿）} & \textbf{Sussman 再初期化} \\
\hline
対象 & $\psi \to \phi$ の点ごとの変換 & $\phi$ のみを直接更新 \\
計算式 & $\phi = \varepsilon\ln(\psi/(1-\psi))$（解析的） & $\partial\phi/\partial\tau + \mathrm{sgn}(\phi_0)(|\bnabla\phi|-1)=0$ \\
コスト & $\Ord{N}$，1 回の対数計算 & $\Ord{N \cdot N_\tau}$，PDE を反復求解 \\
精度 & $\phi = H_\varepsilon^{-1}(\psi)$ の厳密な逆変換 & $|\bnabla\phi|=1$ を修正するが $\psi$ との整合性なし \\
適合 & CLS 法（$\psi$ が基本変数） & 標準 LS 法（$\phi$ が基本変数） \\
\hline
\end{tabular}
\end{center}

\textbf{本稿の選択：} CLS 法では $\psi$ を移流して保存性を得るのが主目的であり，
$\phi$ はあくまで幾何量（法線・曲率）計算の補助として使う．
ロジット関数は各格子点で $\Ord{N}$ の計算（対数1回）で完結し，
$\psi$ と $\phi$ の整合性を代数的に厳密に保てるため，CLS 法に自然に適合する．
飽和域（$\psi \to 0$ または $\psi \to 1$）のみニュートン法のフォールバックを用いるが，
対象格子点数は $\Ord{e^{-15}}$ 以下であり計算コストは無視できる（付録~\ref{app:newton_quad_conv} 参照）．
Sussman 再初期化は $\phi$ の符号付き距離性質（$|\bnabla\phi|=1$）を直接維持したい
標準 LS 法に適しており，CLS 法との組み合わせでは過剰なコストと $\psi$--$\phi$ 不整合を生む恐れがある．
\end{tcolorbox}

% ═══════════════════════════════════════════════════════════════════════════════
\subsection{\texorpdfstring{$\psi$}{ψ} 微分による曲率の直接計算：ロジット逆変換の省略}
\label{sec:curvature_direct_psi_cls}
% ═══════════════════════════════════════════════════════════════════════════════

第~\ref{sec:governing}章の定理~\ref{thm:curvature_invariance} により，
曲率は $\phi$ と $\psi$ のいずれから計算しても同一であることが示された．
本節では，この不変性を CLS 法に適用し，
ロジット逆変換（§~\ref{sec:newton_inversion}）を\textbf{省略}して
$\psi$ の微分値から曲率を直接計算する具体的手法を述べる．

\begin{tcolorbox}[resultbox, title={CLS 法における $\psi$ 微分による曲率の直接計算}]
\label{result:curvature_direct_psi}
式~\eqref{eq:curvature_2d} において $\phi$ を $\psi$ で置き換えた
\begin{equation}
  \kappa = -\frac{
    \psi_y^2\,\psi_{xx}
    - 2\psi_x\psi_y\,\psi_{xy}
    + \psi_x^2\,\psi_{yy}
  }{\bigl(\psi_x^2 + \psi_y^2\bigr)^{3/2}}
  \label{eq:curvature_psi_2d}
\end{equation}
は数学的に厳密に式~\eqref{eq:curvature_2d} と同一の曲率を与える
（定理~\ref{thm:curvature_invariance}）．
CCD 法は $\psi$ の解法過程で $\psi$, $\psi'$, $\psi''$ を
同時に $\Ord{h^6}$ 精度で算出するため，
式~\eqref{eq:curvature_psi_2d} の全ての入力が追加計算なしに利用可能である．

\medskip
\noindent\textbf{従来法（$\psi \to \phi \to \kappa$）に対する利点：}
\begin{enumerate}
  \item ロジット逆変換（式~\eqref{eq:logit_inverse}）の\textbf{完全な省略}：
        飽和領域（$\psi \to 0, 1$）での対数関数の特異性による誤差増幅が
        原理的に排除される．
  \item 逆変換後の $\phi$ に対する差分演算が不要となり，
        $\psi$ の滑らかな関数構造がそのまま保存される．
  \item 再初期化による Eikonal 条件 $|\bnabla\phi| \approx 1$ の維持が
        曲率精度の前提条件ではなくなる：
        曲率はレベルセットの幾何学的性質であり，$|\bnabla\phi|$ の値には依存しない．
\end{enumerate}
\end{tcolorbox}

% ─────────────────────────────────────────────────────────────────────────────
\subsubsection{\texorpdfstring{数値的安定性：$|\bnabla\psi|$ の挙動と適用領域}{数値的安定性：|∇ψ| の挙動と適用領域}}
% ─────────────────────────────────────────────────────────────────────────────

式~\eqref{eq:curvature_psi_2d} の分母
$(\psi_x^2+\psi_y^2)^{3/2}$ は，$\psi$ のプロファイル形状に依存する．
界面近傍（$\psi \approx 0.5$）ではシグモイド勾配が最大値
$|\bnabla\psi|_{\max} = 1/(4\varepsilon)$（$|\bnabla\phi|=1$ のとき）
をとり，安定な除算が可能である．
一方，$\psi \to 0$ または $\psi \to 1$ の領域では
$|\bnabla\psi| \to 0$ となるが，
これらの領域は界面から遠く，表面張力項の評価に曲率が寄与しない．
したがって，実装上は以下のハイブリッド戦略をとる：
\begin{itemize}
  \item \textbf{界面近傍}（$\psi_{\min} < \psi < 1-\psi_{\min}$，
        推奨 $\psi_{\min}=0.01$）：
        式~\eqref{eq:curvature_psi_2d} を適用し，高精度な曲率を直接計算する．
  \item \textbf{界面遠方}（$\psi \le \psi_{\min}$ または
        $\psi \ge 1-\psi_{\min}$）：
        $\kappa = 0$ とする（表面張力寄与なし）．
\end{itemize}
この閾値選択に対するロバスト性は高い：
表面張力項 $\kappa\,\bnabla\psi/We$ において
$|\bnabla\psi| \to 0$ の領域では $\kappa$ が不定となっても
積 $\kappa\,|\bnabla\psi|$ は有界に留まるためである．

% ─────────────────────────────────────────────────────────────────────────────
\subsubsection{実効界面厚み \texorpdfstring{$\varepsilon_{\textup{eff}}$}{ε\_eff} によるプロファイル品質監視}
% ─────────────────────────────────────────────────────────────────────────────

$\psi$ による直接曲率計算では $\varepsilon$ の値は不要であるが，
数値拡散や再初期化の不完全さによるプロファイル歪みを
\textbf{診断}する目的で，各格子点における実効界面厚みを定義できる：
\begin{equation}
  \varepsilon_{\textup{eff},i}
  = \frac{\psi_i(1-\psi_i)}{|\bnabla\psi|_i}
  \label{eq:epsilon_eff}
\end{equation}
この定義は，シグモイド微分公式
$|\bnabla\psi| = \frac{1}{\varepsilon}\psi(1-\psi)\,|\bnabla\phi|$
と Eikonal 条件 $|\bnabla\phi| \approx 1$ から導かれる．
界面近傍で $\varepsilon_{\textup{eff}} \approx \varepsilon$（設計値）であれば
プロファイルが正常に維持されていることを意味し，
$\varepsilon_{\textup{eff}} \gg \varepsilon$ は
数値拡散による過度な界面幅増大を示す診断指標となる．

% ═══════════════════════════════════════════════════════════════════════════════
\subsection{曲率の数値計算と安定化：再初期化との接続}
\label{sec:curvature_stab}
% ═══════════════════════════════════════════════════════════════════════════════

第~\ref{sec:governing}章 §~\ref{sec:curvature}（式~\eqref{eq:curvature_2d}）で曲率の展開式を導出し，
§~\ref{sec:csf} の box~\ref{box:curvature_stab} で実装クイックスタート（分母クランプ式~\eqref{eq:curvature_clamped}）を提示した．
\textbf{本節の目的}は，そのクランプが「なぜ必要か」を CLS 再初期化の文脈から説明することである：
ゼロ除算は再初期化の不完全性に起因するため，クランプは応急処置であり，
根本的な保証は再初期化の品質にある——この因果関係を以下で定量的に示す．
なお，前節の $\psi$ 直接法（式~\eqref{eq:curvature_psi_2d}）を用いる場合は
Eikonal 条件への依存が解消されるため，本節の正則化は不要となる．

第~\ref{sec:governing}章（式~\eqref{eq:curvature_2d}）で導出した2次元曲率の展開式
\[
  \kappa = -\frac{
    \phi_y^2\,\phi_{xx} - 2\phi_x\phi_y\,\phi_{xy} + \phi_x^2\,\phi_{yy}
  }{\bigl(\phi_x^2 + \phi_y^2\bigr)^{3/2}}
\]
は数学的には厳密であるが，数値計算で直接使用するには注意が必要である．
問題の核心は\textbf{分母 $(\phi_x^2+\phi_y^2)^{3/2}$ のゼロ除算リスク}であり，
これは CLS 再初期化の品質（Eikonal 条件 $|\bnabla\phi|=1$ の維持）と
密接に関連した実装上の重要課題である——以下に詳述する．

% ─────────────────────────────────────────────────────────────────────────────
\subsubsection{ゼロ除算の発生メカニズム：再初期化の不完全性との関係}
% ─────────────────────────────────────────────────────────────────────────────

真の符号付き距離関数 $\phi$ であれば，Eikonal 条件 $|\bnabla\phi|=1$ が
至るところで成立するため，分母は常に $1^{3/2}=1$ となりゼロ除算は起きない．
しかし，数値計算では再初期化の不完全性により $|\bnabla\phi|$ が局所的に小さくなることがある．
この場合，式~\eqref{eq:curvature_2d} の分母 $(\phi_x^2+\phi_y^2)^{3/2}$ がゼロに近づき，
$\kappa$ が非物理的な大値をとって表面張力項が発散する．

% ─────────────────────────────────────────────────────────────────────────────
\subsubsection{ゼロ除算の回避：分母への下限値の導入}
% ─────────────────────────────────────────────────────────────────────────────

\paragraph{ゼロ除算の回避（必須実装）}
式~\eqref{eq:curvature_2d} の数値実装では，分母にゼロ除算防止のための
\textbf{下限値 $\varepsilon_{\text{norm}}^2$} を加える（式~\eqref{eq:curvature_clamped}，§~\ref{sec:curvature} 参照）．
$|\bnabla\phi|^2 \gg \varepsilon_{\text{norm}}^2$ の正常な界面近傍では元の式と同一の結果を与え，
$|\bnabla\phi|^2 \lesssim \varepsilon_{\text{norm}}^2$ の特異点付近では
分母が $\varepsilon_{\text{norm}}^3$ に固定されて $\kappa$ が有界に保たれる．
再初期化が正常なら $|\bnabla\phi| \approx 1$ であるため，
$\varepsilon_{\text{norm}} = 10^{-3} \sim 10^{-4}$（次元なし）で十分である．
$\varepsilon_{\text{norm}}$ が大きすぎると曲率の系統的過小評価が生じ，
相対誤差は $\Ord{\varepsilon_{\text{norm}}^2/|\bnabla\phi|^2}$ 程度となる．
\textbf{推奨値：} $\varepsilon_{\text{norm}} = 10^{-3}$．
密度比 $\rho_l/\rho_g \gg 1$ で Eikonal 誤差が大きい場合は $10^{-2}$ も検討する．

$|\bnabla\phi|^2 = 1 + \delta^2$（$\delta$ は Eikonal 誤差）の場合，$|\bnabla\phi|^2 \gg \varepsilon_{\text{norm}}^2$ のとき $\max$ は $|\bnabla\phi|^2$ を選ぶので：
\begin{align*}
  \kappa_{\text{reg}} &= -\frac{\text{分子}}{(1+\delta^2)^{3/2}}
  \approx -\frac{\text{分子}}{1} \cdot \left(1 - \frac{3}{2}\delta^2 + \cdots\right)
\end{align*}
Eikonal 誤差 $\delta \ll 1$ のとき曲率への影響は $\Ord{\delta^2}$ に留まり，
再初期化が概ね正常に機能している限り高精度が維持される．
すなわち，\textbf{曲率精度の第一の保証は正則化ではなく，再初期化の品質（Eikonal 条件の維持）にある}．

% ─────────────────────────────────────────────────────────────────────────────
\subsubsection{曲率計算の実装チェックリスト}
% ─────────────────────────────────────────────────────────────────────────────

\begin{tcolorbox}[algbox, title={曲率計算の実装チェックリスト}]

\textbf{推奨経路（$\psi$ 直接法）}——定理~\ref{thm:curvature_invariance}（§~\ref{sec:curvature_direct_psi}）に基づき，
ロジット逆変換を省略して $\psi$ の微分値から曲率を直接計算する：

\begin{enumerate}
  \item \textbf{1次・2次微分の計算：}
        移流後の $\psi^{n+1}$ に対して
        $\psi_x, \psi_y, \psi_{xx}, \psi_{yy}, \psi_{xy}$ を CCD 法（第~\ref{sec:CCD}章）で
        $\Ord{h^6}$ 精度で計算する（CCD は解法過程で同時に算出）．

  \item \textbf{適用領域の判定：}
        界面近傍（$\psi_{\min} < \psi < 1-\psi_{\min}$，推奨 $\psi_{\min}=0.01$）の
        格子点のみを曲率計算の対象とし，界面遠方では $\kappa = 0$ とする．

  \item \textbf{曲率の計算：}
        式~\eqref{eq:curvature_psi_2d} で $\kappa$ を計算する．

  \item \textbf{（オプション）プロファイル品質監視：}
        実効界面厚み $\varepsilon_{\textup{eff}}$（式~\eqref{eq:epsilon_eff}）を
        評価し，$\varepsilon_{\textup{eff}}/\varepsilon$ が設計値から大きく逸脱していないか確認する．

  \item \textbf{（オプション）ガウシアンフィルタ：}
        $\kappa$ に $3\times 3$ ガウシアンフィルタを適用して
        高周波数ノイズを低減する．高密度比問題では特に有効．

  \item \textbf{単体テスト（半径 $R$ の円）：}
        解析値 $\kappa = -1/R$ との誤差 $E_\kappa(h) = O(h^6)$ 収束を
        $N = 16, 32, 64, 128$ で確認する（§~\ref{sec:curvature_circle_example}）．
\end{enumerate}

\medskip
\textbf{代替経路（$\phi$ 経由）}——Extension PDE（§~\ref{sec:extension_pde}）で $\phi$ が必要な場合：
\begin{enumerate}
  \item ロジット逆変換（§~\ref{sec:newton_inversion}）で $\phi^{n+1}$ を取得．
  \item CCD 法で $\phi$ の1次微分を計算し，界面法線 $\hat{\bm{n}} = \bnabla\phi/|\bnabla\phi|$ を求める．
  \item Extension PDE $\partial q/\partial\tau + S(\phi)\,\hat{\bm{n}}\cdot\bnabla q = 0$
        により，圧力場 $\delta p$ を界面越しに滑らかに延長する．
  \item 延長後の滑らかな場に対して CCD $\bnabla(\delta p)$ を適用し，$\Ord{h^6}$ 精度の速度補正を得る．
\end{enumerate}
\end{tcolorbox}

% ═══════════════════════════════════════════════════════════════════════════════
\subsection{Extension PDE による界面越し場延長：CCD 高次精度の維持}
\label{sec:extension_pde}
% ═══════════════════════════════════════════════════════════════════════════════

§~\ref{sec:curvature_direct_psi_cls} で述べた $\psi$ 直接法により，
曲率 $\kappa$ は界面近傍まで $\Ord{h^6}$ 精度で計算できる．
しかし，CSF 表面張力（§~\ref{sec:csf}）による\textbf{圧力場 $p$} には
界面を跨ぐ Laplace ジャンプ $\Delta p = \kappa/\mathrm{We}$ が存在し，
この不連続性を CCD ステンシル $(i{-}2,\ldots,i{+}2)$ が参照すると
Gibbs 振動が速度補正 $\bnabla(\delta p)$ に混入する．

\textbf{Extension PDE}~\cite{Aslam2004} は，圧力場を界面越しに
法線方向へ定数外挿することで，CCD が参照する場を $C^\infty$ に保つ：
\begin{equation}
  \pd{q}{\tau} + S(\phi)\,\hat{\bm{n}}\cdot\bnabla q = 0
  \label{eq:extension_pde}
\end{equation}
ここで $q$ は延長対象の物理量（$\delta p$ または $p^n$），
$S(\phi) = \mathrm{sign}(\phi)$ は伝搬方向の符号，
$\hat{\bm{n}} = \bnabla\phi/|\bnabla\phi|$ は界面法線である．
界面法線 $\hat{\bm{n}}$ は $\phi$（滑らかな符号付き距離関数）から
CCD $D^{(1)}$ で $\Ord{h^6}$ 精度に計算される．

式~\eqref{eq:extension_pde} を仮想時間 $\tau$ 方向に解くと，
各相の値が法線方向に一定のまま対岸へ伝搬し，
\textbf{界面の両側で $q$ が滑らか}となる（図~\ref{fig:extension_schematic}）．

\begin{figure}[htbp]
\centering
\begin{tikzpicture}[scale=0.85]
  % 延長前
  \begin{scope}[xshift=-50mm]
    \node[font=\small\bfseries, text=navy] at (0,2.8) {延長前：$\delta p$ に不連続};
    \draw[->] (-2.5,0) -- (2.5,0) node[right,font=\small]{$x$};
    \draw[->] (0,-0.5) -- (0,2.5) node[above,font=\small]{$\delta p$};
    \draw[very thick, blue2, domain=-2.2:0] plot (\x, {1.8});
    \draw[very thick, red2, domain=0:2.2] plot (\x, {0.3});
    \draw[dashed, gray] (0,-0.3) -- (0,2.3);
    \node[font=\scriptsize] at (0,-0.5) {$\Gamma$};
    \node[font=\scriptsize, blue2] at (-1.2, 2.1) {液体};
    \node[font=\scriptsize, red2] at (1.2, 0.6) {気体};
    \draw[thick, orange2, decorate, decoration={zigzag, amplitude=3pt, segment length=5pt}]
      (-0.3, 1.0) -- (0.3, 1.0);
    \node[font=\scriptsize, orange2] at (0, 0.6) {Gibbs};
  \end{scope}
  % 矢印
  \draw[-Stealth, very thick, navy] (-1.5, 1.2) -- (1.0, 1.2)
    node[midway, above, font=\small, navy]{Extension PDE};
  % 延長後
  \begin{scope}[xshift=50mm]
    \node[font=\small\bfseries, text=navy] at (0,2.8) {延長後：$\delta p_\mathrm{ext}$ は滑らか};
    \draw[->] (-2.5,0) -- (2.5,0) node[right,font=\small]{$x$};
    \draw[->] (0,-0.5) -- (0,2.5) node[above,font=\small]{$\delta p_\mathrm{ext}$};
    \draw[very thick, blue2, domain=-2.2:2.2] plot (\x, {1.8});
    \draw[dashed, gray] (0,-0.3) -- (0,2.3);
    \node[font=\scriptsize] at (0,-0.5) {$\Gamma$};
    \node[font=\scriptsize, blue2] at (1.2, 2.1) {延長済み};
    \node[font=\scriptsize, teal2] at (0, 0.4) {CCD $\bnabla(\delta p_\mathrm{ext})$: $\Ord{h^6}$};
  \end{scope}
\end{tikzpicture}
\caption{Extension PDE による圧力場の界面越し延長（概念図）．
左：CSF 表面張力による Laplace ジャンプが CCD ステンシルを跨ぐと Gibbs 振動が発生．
右：Extension PDE で液体側の値を気体側へ滑らかに延長した後は，
CCD $\bnabla(\delta p)$ が $\Ord{h^6}$ 精度を界面直近まで維持する．}
\label{fig:extension_schematic}
\end{figure}

\subsubsection{離散化と CLS 再初期化との構造的類似}

式~\eqref{eq:extension_pde} は純粋な双曲型 PDE であり，
§~\ref{sec:reinit} の CLS 再初期化方程式~\eqref{eq:cls_reinit_iso} の圧縮項
$\bnabla\cdot[\psi(1{-}\psi)\hat{\bm{n}}]$ と\textbf{同型の構造}を持つ．
離散化は以下の通り：
\begin{itemize}
  \item \textbf{空間微分：} $\hat{\bm{n}}\cdot\bnabla q$ を1次精度風上差分で評価する．
        延長対象 $q$ が界面で不連続であるため，CCD は適用できない
        （延長\textbf{後}の滑らかな場には CCD を安全に適用できる）．
  \item \textbf{時間積分：} 前進 Euler 法．CFL 条件は $\Delta\tau \leq 0.5\,h_{\min}$
        （特性速度 $|S\hat{\bm{n}}| \leq 1$）．
  \item \textbf{反復回数：} CCD ステンシル幅（$2$ セル）を延長でカバーするには
        $n_\mathrm{ext} \approx 4$--$8$ 回で十分（推奨デフォルト $n_\mathrm{ext} = 5$）．
\end{itemize}

\begin{tcolorbox}[resultbox, title={CSF + Extension PDE：CCD 高次精度との整合}]
\label{result:csf_extension}
Extension PDE は界面力モデルではなく，\textbf{場の滑らか化手法}である．
表面張力の扱いは CSF（Balanced-Force，§~\ref{sec:csf}）を維持する：
\[
  \bm{f}_\sigma = \frac{\kappa}{\mathrm{We}}\bnabla\psi
\]
ここで $\kappa$ と $\bnabla\psi$ はともに CCD $\Ord{h^6}$ で計算される．
CSF モデル誤差 $\Ord{\varepsilon^2}$ は依然として全体精度の律速であるが，
Extension PDE により\textbf{それ以外の全コンポーネント}が界面直近まで $\Ord{h^5}$ 以上を維持し，
CSF モデル誤差のみが純粋に測定可能となる．

\medskip
\noindent\textbf{Extension PDE の適用対象（§~\ref{sec:algorithm} Step 5c/6b）：}
\begin{enumerate}
  \item \textbf{IPC 圧力 $p^n$}（Predictor 前）：
        CCD $\bnabla p^n$ の Gibbs 振動を防止．
  \item \textbf{圧力増分 $\delta p$}（速度補正前）：
        CCD $\bnabla(\delta p)$ の界面近傍精度を $\Ord{h^6}$ に維持．
\end{enumerate}
$\psi$, $\rho$, $\mu$ は CLS 正則化 $H_\varepsilon$ で既に滑らかであるため延長不要．
曲率 $\kappa$ は $\psi$ の直接法（式~\eqref{eq:curvature_psi_2d}）で計算するため延長不要．

\medskip
\noindent PPE 求解（§~\ref{sec:gfm}）との関係：本節の Extension PDE が $p^n$ および $\delta p$ を
界面越しに延長することで，§~\ref{sec:ccd_poisson_matrix} の CCD-Poisson 演算子が
界面直近まで $\Ord{h^6}$ 精度を維持できる．
\end{tcolorbox}
